\documentclass[fleql]{beamer}

\usetheme[
  progressbar=frametitle,
  numbering=fraction,
  sectionpage=none,
  subsectionpage=none
]{metropolis}

\title{Week 5 \\ Understanding under Threat}
\subtitle{Quantum Mechanics and Artificial Intelligence}
\author{Hans Halvorson}
\date{May 2026}

\usepackage[backend=biber, style=authortitle]{biblatex}
\addbibresource{refs.bib}
\setlength{\bibitemsep}{0.5em}
\AtBeginBibliography{\scriptsize}

\usepackage{amsmath,amssymb}
\usepackage{soul}
\usepackage{tikz}
\usetikzlibrary{decorations.pathmorphing}

\newcommand{\kw}[1]{\textbf{#1}}
\newcommand{\ket}[1]{|#1\rangle}
\newcommand{\bra}[1]{\langle #1|}
\newcommand{\ip}[2]{\langle #1|#2\rangle}

\setbeamertemplate{section in toc}[sections numbered]
\setbeamerfont{section in toc}{size=\small}

\AtBeginSection[]
{
  \begin{frame}
    \frametitle{Table of Contents}
    \tableofcontents[currentsection,hideallsubsections,part=\thepart]
  \end{frame}
}

\setbeamertemplate{footline}[frame number]
\setbeamertemplate{navigation symbols}{}
\setbeamertemplate{part page}{}

\setcounter{secnumdepth}{1}

\begin{document}

\begin{frame}
  \titlepage
\end{frame}

\begin{frame}{Today's plan}

  \textbf{Part I: Quantum Mechanics} \hfill \emph{(lectures 1--2)}
  \begin{itemize}\small
  \item The framework of QM: states, observables, dynamics
  \item The measurement problem and the Copenhagen myth
  \item Interpretations: GRW, Bohm, Everett, \ldots
  \item Complementarity: Bohr's response to the crisis
  \item Bell's theorem and the limits of causal understanding
  \end{itemize}

  \vspace{0.8em}
  \textbf{Part II: Artificial Intelligence} \hfill \emph{(lecture 3)}
  \begin{itemize}\small
  \item Models and understanding: what philosophy of science says
  \item The double meaning of ``model''
  \item Sullivan's argument: link uncertainty vs.\ implementation opacity
  \item A deeper challenge: is there an inside?
  \end{itemize}

\end{frame}

%%% ============================================================
%%% PART I: QUANTUM MECHANICS
%%% ============================================================

\part{Quantum Mechanics}

\begin{frame}{Group Discussion}

  From last week, quantum mechanics is our recurring example of
  a theory where prediction and understanding seem to come apart.

  \vspace{0.8em}
  \begin{enumerate}
  \item Do you feel that you \textbf{understand} quantum mechanics?
    What would it even mean to understand it?
  \item Is ``shut up and calculate'' a satisfying attitude toward a
    fundamental theory of nature?
  \end{enumerate}

\end{frame}

\section{Introduction}

\begin{frame}{Two perspectives today}

  \textbf{First-order}\enspace Which interpretation of QM is
  correct? What is the theory \emph{actually} saying about
  the world?

  \bigskip

  \textbf{Second-order}\enspace There has been no consensus on
  the foundations of QM for \emph{100 years}. Should we be
  concerned? What does this tell us about science as a
  knowledge-producing enterprise?

  \bigskip

  Both questions bear directly on scientific understanding ---
  which is why De Regt and Dieks chose quantum non-locality as
  their motivating example.

\end{frame}

\begin{frame}{The foundational crisis: 1900--1935}

  \begin{itemize}\small
  \item 1900: Planck introduces the quantum of action $h$
  \item 1905: Einstein explains the photoelectric effect using photons
  \item 1913: Bohr's model of the hydrogen atom
  \item 1925--26: Heisenberg's matrix mechanics; Schr\"{o}dinger's
    wave mechanics
  \item 1927: Born's probability interpretation; Heisenberg's
    uncertainty principle; Bohr's complementarity
  \item 1935: The EPR argument --- QM must be incomplete
  \end{itemize}

  \vspace{0.4em}
  \small The early pioneers were acutely aware that they did
  \emph{not} understand what was going on. Bohr thought our
  ability to understand the physical world was genuinely up
  for grabs.

\end{frame}

\begin{frame}{Apparent resolution: the Copenhagen school}

  \begin{itemize}
  \item By 1927 a working consensus formed around Bohr's school
  \item Bohr's key insight: the existence of the \kw{quantum of
    action} calls for a \emph{revision} of what it means to give an
    objective, unambiguous description of physical phenomena
    \begin{itemize}
    \item An unambiguous description must specify the experimental
      arrangement; the conditions of observation become part of
      the description itself
    \item Classical concepts remain indispensable for describing
      the apparatus, but cannot all be simultaneously applied to
      the quantum object
    \end{itemize}
  \item This leads to \kw{complementarity}: a structural feature of
    unambiguous description, not a limitation on knowledge
  \item This is a positive philosophical proposal, not a
    counsel of despair
  \end{itemize}

\end{frame}

\begin{frame}{Revival of the crisis}

  \begin{itemize}
  \item 1952: Bohm publishes a working hidden-variable
    interpretation --- shattering the apparent consensus
  \item 1957: Everett proposes the relative-state (many-worlds)
    interpretation
  \item 1964: Bell proves his non-locality theorem
  \item Today: active foundational debate; the cutting edge of
    physics proceeds without settling the foundations
  \end{itemize}

  \vspace{0.5em}
  \textbf{The second-order question:} is 100 years without consensus
  acceptable? Does it reveal something deep about QM, or is it a
  sociological failure of physics?

\end{frame}

\section{Framework of Quantum Mechanics}

\begin{frame}{States as vectors}

  \begin{itemize}
  \item The state of a quantum system is a unit vector $\ket{\psi}$
    in a complex Hilbert space $\mathcal{H}$
  \item \kw{Superposition principle}: if $\ket{\psi}$ and $\ket{\phi}$
    are states, so is $\alpha\ket{\psi} + \beta\ket{\phi}$
    (after normalization)
  \item For a \kw{qubit} (two-level system): $\mathcal{H} = \mathbb{C}^2$
  \item Computational basis: $\ket{0} \equiv \ket{\uparrow}$,\quad
    $\ket{1} \equiv \ket{\downarrow}$
  \end{itemize}

\end{frame}

\begin{frame}{The Bloch sphere}

\begin{center}
\begin{tikzpicture}[scale=1.1]
  \draw[thick] (0,0) circle (1.5);
  \draw[gray!50] (1.5,0) arc (0:180:1.5 and 0.45);
  \draw[gray!50,dashed] (-1.5,0) arc (180:360:1.5 and 0.45);
  \draw[->,semithick] (0,-1.9) -- (0,2.1) node[above] {\scriptsize$z$};
  \draw[->,semithick] (-1.9,0) -- (2.2,0) node[right] {\scriptsize$x$};
  \draw[->,semithick,dashed] (0,0) -- (-0.85,-0.6)
    node[below left] {\scriptsize$y$};
  \fill (0,1.5) circle (2pt)
    node[above,font=\scriptsize] {$\ket{0} = \ket{\uparrow}_z$};
  \fill (0,-1.5) circle (2pt)
    node[below,font=\scriptsize] {$\ket{1} = \ket{\downarrow}_z$};
  \fill (1.5,0) circle (1.5pt)
    node[above right,font=\scriptsize] {$\ket{+}_x$};
  \fill (-1.5,0) circle (1.5pt)
    node[above left,font=\scriptsize] {$\ket{-}_x$};
  \draw[->,thick,blue!75] (0,0) -- (0.9,1.1)
    node[right,font=\scriptsize] {$\ket{\psi}$};
  \fill[blue!75] (0,0) circle (1.5pt);
\end{tikzpicture}
\end{center}

\vspace{-0.3em}
\small Every pure state of a qubit is a point on the sphere.
Antipodal points are orthogonal.

\end{frame}

\begin{frame}{Observables as operators}

  \begin{itemize}
  \item Physical quantities: \kw{Hermitian operators} $A = A^\dagger$
  \item Possible measurement outcomes: eigenvalues $\lambda_i$
  \item Eigenstates $\ket{a_i}$: states with definite value $\lambda_i$
  \end{itemize}

  \vspace{0.5em}
  \textbf{Spin operators (Pauli matrices $\times\,\hbar/2$):}
  \[
  S_z = \frac{\hbar}{2}\begin{pmatrix}1&0\\0&-1\end{pmatrix},\quad
  S_x = \frac{\hbar}{2}\begin{pmatrix}0&1\\1&0\end{pmatrix},\quad
  S_y = \frac{\hbar}{2}\begin{pmatrix}0&-i\\i&0\end{pmatrix}
  \]

\end{frame}

\begin{frame}{Superpositions and the puzzle}

  If the system is in $\ket{\psi} = \alpha\ket{a_1}+\beta\ket{a_2}$,
  observable $A$ has \emph{no definite value} --- yet measurement
  always yields a definite outcome $\lambda_1$ or $\lambda_2$.

  \vspace{0.5em}
  \textbf{The puzzle:} what is the system doing between preparation
  and measurement? Does the question even make sense?

  \vspace{0.5em}
  This tension is the engine of every interpretation of QM.

\end{frame}

\begin{frame}{Uncertainty relations}

  Non-commuting observables cannot both have definite values:
  \[
  [S_y,\,S_z] = i\hbar S_x
  \quad\Longrightarrow\quad
  \Delta S_y \cdot \Delta S_z \;\geq\;
  \tfrac{\hbar}{2}|\langle S_x\rangle|
  \]

  \textbf{Spin-$z$ eigenstate decomposed in spin-$y$ basis:}
  \[
  \ket{\uparrow}_z \;=\;
  \tfrac{1}{\sqrt{2}}\ket{\uparrow}_y + \tfrac{1}{\sqrt{2}}\ket{\downarrow}_y
  \]
  so $\Delta S_y = \hbar/2$ is maximal when $S_z$ is definite.

  \vspace{0.3em}
  Conceptually: \kw{Fourier analysis} --- narrow in one basis,
  spread in the conjugate basis.

\end{frame}

\begin{frame}{Composite systems and entanglement}

  Composite systems: $\mathcal{H}_{AB} =
  \mathcal{H}_A \otimes \mathcal{H}_B$, with
  $x\otimes(y+z) = (x\otimes y)+(x\otimes z)$.

  \vspace{0.3em}
  Not all states are product states. The \kw{singlet state}
  \[
  \ket{\Psi^-} = \tfrac{1}{\sqrt{2}}
  \bigl(\ket{\uparrow\downarrow} - \ket{\downarrow\uparrow}\bigr)
  \]
  cannot be written $\ket{\psi}_A\otimes\ket{\phi}_B$:
  it is \kw{entangled}.

  \vspace{0.5em}
  \textbf{Discussion (2 min):} what features does a system have if
  it is entangled with another? What does it mean for its ``state''?

\end{frame}

\begin{frame}{Linear dynamics and the Born rule}

  Time evolution:
  $i\hbar\frac{d}{dt}\ket{\psi} = H\ket{\psi}$,\quad
  solution $\ket{\psi(t)} = U(t)\ket{\psi(0)}$ (\kw{unitary}).

  \vspace{0.3em}
  Crucially, $U$ is \kw{linear}:
  $U(\alpha\ket{\psi}+\beta\ket{\phi}) = \alpha U\ket{\psi}+\beta U\ket{\phi}$.
  This is exactly what generates the measurement problem.

  \vspace{0.5em}
  \textbf{Born rule:}
  \[
  \mathrm{Pr}_\psi(A=\lambda) =
  \langle\psi,\,E^A(\lambda)\,\psi\rangle
  \]
  For $\ket{\psi}=\sum_i c_i\ket{a_i}$:\quad
  $\mathrm{Pr}_\psi(A=\lambda_i) = |c_i|^2$

  \vspace{0.4em}
  \textbf{Discussion:} do we \emph{understand} QM if it is
  just a recipe for computing these probabilities?

\end{frame}

\section{The Measurement Problem}

\begin{frame}{Definition and derivation}

  \textbf{Ideal measurement of $A$:}
  $\ket{a_i}\otimes\ket{r} \;\longrightarrow\; \ket{a_i}\otimes\ket{y_i}$
  (each eigenstate paired with a distinct pointer state --- unitary).

  \vspace{0.5em}
  Now suppose $\ket{\psi} = \sum_i c_i\ket{a_i}$.
  By \textbf{linearity}:
  \[
  \Bigl(\sum_i c_i\ket{a_i}\Bigr)\otimes\ket{r}
  \;\longrightarrow\;
  \sum_i c_i\,\bigl(\ket{a_i}\otimes\ket{y_i}\bigr)
  \]

  The meter ends in an \textbf{entangled superposition} --- no
  definite reading. But we \emph{always} observe a definite outcome.
  \textbf{Contradiction.}

\end{frame}

\begin{frame}{The collapse postulate}

  \textbf{Standard ``solution'':}
  $\sum_i c_i\ket{a_i} \longrightarrow \ket{a_k}$
  with probability $|c_k|^2$.

  \vspace{0.5em}
  \textbf{Problems:}
  \begin{itemize}
  \item Violates unitarity and linearity --- the two most fundamental
    features of the theory
  \item Creates a two-dynamics theory:
    Schr\"{o}dinger evolution \emph{and} collapse
  \item When does collapse occur? Where is the ``Heisenberg cut''?
  \item The theory gives no definition of ``measurement''
  \end{itemize}

\end{frame}

\begin{frame}{The Copenhagen myth}

  \begin{itemize}
  \item Von Neumann (1932) first axiomatized collapse in
    \emph{Grundlagen der Quantenmechanik}; by the 1950s it was
    standard textbook orthodoxy
  \item The label \kw{Copenhagen interpretation} was applied
    retroactively --- Bohr never used it
  \item In reality there was no single Copenhagen view: Bohr,
    Heisenberg, Born, Pauli all held different positions
  \item \textbf{HH:} Bohr never claimed the measurement problem is
    \emph{solved} by collapse --- his actual response is
    philosophically much deeper, and has been systematically misread
  \end{itemize}

\end{frame}

\begin{frame}{Decoherence: not a solution}

  \kw{Decoherence}: environmental interaction entangles the system
  rapidly; the \emph{reduced} state approaches a mixture.

  \vspace{0.5em}
  \textbf{Why this does not solve the measurement problem:}
  \begin{enumerate}
  \item The reduced state is not \emph{uniquely} interpretable as
    a mixture over definite pointer readings
  \item The global state (system $+$ environment) is still a
    superposition --- decoherence is a local approximation,
    not a collapse
  \end{enumerate}

  \vspace{0.3em}
  \small Decoherence explains why interference is hard to observe
  macroscopically. It does not explain why we observe \emph{one}
  outcome.

\end{frame}

\section{Interpretations of Quantum Mechanics}

\begin{frame}{Four strategies}

  \begin{enumerate}
  \item \kw{Reject unitary dynamics}: add stochastic collapse
    \emph{(GRW)}
  \item \kw{Add hidden variables}: supplement QM with additional
    ontology \emph{(Bohm)}
  \item \kw{Reject uniqueness of outcomes}: all branches are real
    \emph{(Everett)}
  \item \kw{Revise the ontological picture}: resist the demand for
    a classical reality behind QM
    \emph{(relational, QBism, pragmatic)}
  \end{enumerate}

\end{frame}

\begin{frame}{GRW, Bohm, and their problems}

  \textbf{GRW} (Ghirardi, Rimini, Weber 1986): add random
  spontaneous localization events. Free parameters tuned so
  macroscopic objects localize (no pointer superpositions) while
  microscopic systems remain unitary. Open: ontology? Relativistic
  extension?

  \vspace{0.5em}
  \textbf{Bohm (1952)}: both wavefunction and particle positions
  are real; the wave guides particles. Particles always have
  definite positions.
  \begin{itemize}\small
  \item Guidance equation requires a preferred frame ---
    inconsistent with special relativity
  \item Violates the spirit of momentum conservation in
    relativistic interactions
  \end{itemize}

\end{frame}

\begin{frame}{Everett / many worlds}

  \textbf{Everett (1957)}: take unitary dynamics completely seriously;
  no collapse; all branches of the wavefunction are equally real.

  \vspace{0.5em}
  We observe one branch because we ourselves are part of the
  wavefunction --- each measurement splits the world.

  \vspace{0.5em}
  \textbf{Problems:}
  \begin{itemize}\small
  \item \emph{Probability:} the theory is deterministic --- where do
    the Born-rule probabilities come from?
  \item \emph{Preferred basis:} what selects the ``world'' branches?
  \item No-hidden-variables theorems imply the branches cannot
    themselves be classical
  \end{itemize}

  \vspace{0.4em}
  \small \textbf{HH:} these approaches are philosophically
  sophisticated but may relocate rather than dissolve the problem.

\end{frame}

\begin{frame}{Less ontological interpretations}

  \textbf{Relational QM} (Rovelli 1996): quantum states are always
  \emph{relative to an observer} --- there is no
  observer-independent quantum state.

  \vspace{0.3em}
  \small The inspiration is special relativity: velocity has no
  absolute value, only values relative to a reference frame.
  Rovelli proposes the same for quantum states. When Alice measures
  a particle and gets a definite outcome, that outcome is real
  \emph{relative to Alice}. The particle may still be in
  superposition relative to Bob, who has not interacted with it.
  Both descriptions are equally valid. The measurement problem is
  \emph{dissolved}: collapse is real, but only relative to an
  observer, never absolute.

  \vspace{0.4em}
  \normalsize \textbf{QBism} (Fuchs--Schack): quantum states encode
  an agent's \emph{beliefs} about future experiences --- the
  wavefunction is epistemic, not ontic.

  \vspace{0.3em}
  \textbf{Pragmatic} (Healey): QM provides a tool for updating
  credences, not a description of a quantum reality.

\end{frame}

\section{Complementarity}

\begin{frame}{Bohr's single-slit thought experiment}

\begin{columns}

\column{0.49\textwidth}
\centering
\textbf{A: bolted screen}
\begin{center}
\begin{tikzpicture}[scale=0.63]
  \fill[gray!55] (0,0.35) rectangle (0.45,2.1);
  \fill[gray!55] (0,-2.1) rectangle (0.45,-0.35);
  \draw[ultra thick] (0.22,2.1) -- (0.22,2.45);
  \draw[ultra thick] (0.22,-2.1) -- (0.22,-2.45);
  \draw[->,thick] (-1.8,0) -- (0,0);
  \draw[->,blue!80,thick] (0.45,0) -- (2.3,0.75);
  \draw[->,blue!50] (0.45,0) -- (2.3,0.15);
  \draw[->,blue!25] (0.45,0) -- (2.3,-0.35);
  \fill[gray!25] (2.3,-2.1) rectangle (2.6,2.1);
  \foreach \y in {0.75,0.15,-0.35} {
    \fill[blue!50] (2.6,\y-0.09) rectangle (3.05,\y+0.09);
  }
\end{tikzpicture}
\end{center}
\small position known $\Rightarrow$ wave picture $\Rightarrow$ \textbf{interference}

\column{0.49\textwidth}
\centering
\textbf{B: free screen}
\begin{center}
\begin{tikzpicture}[scale=0.63]
  \fill[gray!55] (0,0.35) rectangle (0.45,2.1);
  \fill[gray!55] (0,-2.1) rectangle (0.45,-0.35);
  \draw[thick,decorate,
    decoration={zigzag,segment length=4pt,amplitude=2pt}]
    (0.22,-2.1) -- (0.22,-2.7);
  \draw[->,thick] (-1.8,0) -- (0,0);
  \draw[->,blue!80,thick] (0.45,0) -- (2.3,0.45);
  \draw[->,red,thick] (0.22,-1.7) -- (0.22,-2.2)
    node[right,font=\tiny] {recoil};
  \fill[gray!25] (2.3,-2.1) rectangle (2.6,2.1);
  \fill[gray!45] (2.6,-0.1) rectangle (3.05,0.55);
\end{tikzpicture}
\end{center}
\small momentum known $\Rightarrow$ particle picture $\Rightarrow$ \textbf{no interference}

\end{columns}

\end{frame}

\begin{frame}{Complementarity and visualizability}

  Two arrangements are \kw{complementary} when they are mutually
  exclusive and each exhaustive within its domain.

  \vspace{0.5em}
  \textbf{Spacetime description}\enspace position, trajectory,
  ``which path'' --- \kw{visualizable} (anskuelig) picture

  \medskip

  \textbf{Causal description}\enspace momentum, energy,
  conservation laws --- dynamical picture

  \vspace{0.5em}
  ``The viewpoint of complementarity presents itself as a rational
  generalization of the very ideal of causality.'' (Bohr 1948)

  \vspace{0.3em}
  \small This connects directly to De Regt \& Dieks: Bohr is
  \emph{not} saying we cannot understand quantum phenomena. He
  is proposing a new kind of understanding that does not require
  a single classical picture.

\end{frame}

\begin{frame}{Bohr's deeper point}

  \begin{itemize}
  \item The quantum of action requires a revision of what counts
    as an \kw{objective description}:
    \begin{itemize}
    \item Objectivity does not disappear --- but it now requires
      specifying the experimental conditions as part of any
      unambiguous statement about a quantum phenomenon
    \item We cannot separate ``how Nature is'' from ``the
      experimental arrangement in which it is observed''
    \end{itemize}
  \item This is not about abandoning objectivity but about
    \emph{reconceiving} it for a domain where the interaction
    between object and instrument cannot be neglected or controlled
  \item The misreading of Bohr as ``shut up and calculate'' has
    distorted the foundations debate for 70 years
  \item \textbf{HH:} Bohr is working in a tradition running
    from Kant through H\o ffding: the conditions of unambiguous
    description are themselves a philosophical achievement
  \end{itemize}

\end{frame}

\section{No-Hidden-Variables Theorems}

\begin{frame}{EPR, von Neumann, and Kochen--Specker}

  \textbf{EPR (1935):} QM is incomplete --- quantum correlations
  must be explained by predetermined hidden variables $\lambda$.

  \vspace{0.5em}
  \textbf{Von Neumann (1932):} no hidden-variable theory can
  reproduce QM given additivity $\mathbb{E}_\lambda(A+B) =
  \mathbb{E}_\lambda(A)+\mathbb{E}_\lambda(B)$ for non-commuting
  $A,B$. Bell (1966): this assumption is physically unmotivated.

  \vspace{0.5em}
  \textbf{Kochen--Specker (1967):} no assignment
  $v(A)\in\mathrm{spec}(A)$ to all observables can satisfy
  functional consistency $v(f(A))=f(v(A))$ and the sum rule for
  commuting pairs. \kw{Value definiteness} is impossible.

  \vspace{0.3em}
  \small Kochen--Specker is more compelling than von Neumann: the
  assumptions are operationally motivated.

\end{frame}

\section{Bell's Theorem}

\begin{frame}{The CHSH scenario}

\begin{center}
\begin{tikzpicture}[scale=0.92]
  \draw[fill=gray!20,thick,rounded corners=3pt]
    (-0.45,-0.4) rectangle (0.45,0.4);
  \node[font=\scriptsize] at (0,0) {source};
  \draw[fill=blue!10,thick,rounded corners=3pt]
    (-4.2,-0.45) rectangle (-2.3,0.45);
  \node[font=\scriptsize] at (-3.25,0) {\textbf{Alice}};
  \node[font=\tiny] at (-3.25,-0.65) {setting: $a$ or $a'$};
  \node[font=\tiny] at (-3.25,0.65) {outcome: $\pm 1$};
  \draw[fill=red!10,thick,rounded corners=3pt]
    (2.3,-0.45) rectangle (4.2,0.45);
  \node[font=\scriptsize] at (3.25,0) {\textbf{Bob}};
  \node[font=\tiny] at (3.25,-0.65) {setting: $b$ or $b'$};
  \node[font=\tiny] at (3.25,0.65) {outcome: $\pm 1$};
  \draw[->,thick] (-0.45,0) -- (-2.3,0);
  \draw[->,thick] (0.45,0) -- (2.3,0);
  \draw[<->,dashed,gray!70] (-2.3,-0.9) -- (2.3,-0.9)
    node[midway,below,font=\tiny,gray!70] {no communication};
\end{tikzpicture}
\end{center}

\vspace{0.3em}
\small The entangled pair is measured far apart. Each party
independently chooses a setting and records a $\pm 1$ outcome.

\end{frame}

\begin{frame}{Local hidden variables and the CHSH inequality}

  Suppose outcomes are determined by $\lambda\sim\rho(\lambda)$,
  with \kw{reality} ($A,B$ definite) and \kw{locality}
  (Alice independent of Bob's setting). Define the
  \kw{correlation} between outcomes:
  \[
  E(a,b) = \int A(a,\lambda)\,B(b,\lambda)\,\rho(\lambda)\,d\lambda
  \]
  i.e.\ the expectation value of the product of Alice's and Bob's
  outcomes ($\pm 1$ each) for settings $a$ and $b$.
  Perfect correlation: $E=+1$; perfect anti-correlation: $E=-1$;
  no correlation: $E=0$.

  Under local hidden variables:
  \[
  \bigl|E(a,b) - E(a,b') + E(a',b) + E(a',b')\bigr| \;\leq\; 2
  \]

  \small \textit{Proof sketch:} for each $\lambda$, the combination
  $A(a)[B(b)-B(b')] + A(a')[B(b)+B(b')]$ takes values in
  $\{-2,0,+2\}$.

\end{frame}

\begin{frame}{Quantum violation}

  For $\ket{\Psi^-} = \tfrac{1}{\sqrt{2}}(\ket{\uparrow\downarrow}
  -\ket{\downarrow\uparrow})$:\quad
  $E(a,b) = -\cos(\theta_a-\theta_b)$

  \vspace{0.3em}
  Choose $a=0^\circ$, $a'=90^\circ$, $b=45^\circ$, $b'=135^\circ$:

  \begin{align*}
  E(a,b)   &= -\cos(45^\circ)  = -\tfrac{1}{\sqrt{2}} \\[3pt]
  E(a,b')  &= -\cos(135^\circ) = +\tfrac{1}{\sqrt{2}} \\[3pt]
  E(a',b)  &= -\cos(45^\circ)  = -\tfrac{1}{\sqrt{2}} \\[3pt]
  E(a',b') &= -\cos(45^\circ)  = -\tfrac{1}{\sqrt{2}}
  \end{align*}

  \begin{align*}
  S &= -\tfrac{1}{\sqrt{2}} - \tfrac{1}{\sqrt{2}}
     - \tfrac{1}{\sqrt{2}} - \tfrac{1}{\sqrt{2}}
   = -\tfrac{4}{\sqrt{2}} = -2\sqrt{2} \\[6pt]
  &\Rightarrow\quad \boxed{|S| = 2\sqrt{2} \approx 2.83 > 2}
  \end{align*}

\end{frame}

\begin{frame}{The Jarrett analysis and what Bell means}

  Jarrett (1984): local hidden variables $=$ two conditions:
  \textbf{PI (parameter independence)}\enspace Alice's outcome
  does not depend on Bob's \emph{setting}

  \smallskip

  \textbf{OI (outcome independence)}\enspace Alice's outcome
  does not depend on Bob's \emph{outcome}

  Bell violation $\Leftrightarrow$ at least one of PI or OI fails.
  \begin{itemize}
  \item \textbf{Orthodox (collapse):} OI fails --- collapse makes
    Alice's probability depend on Bob's \emph{outcome}
  \item \textbf{Bohm:} PI fails --- the quantum potential makes
    Alice's probability depend on Bob's \emph{setting}
  \item \textbf{Everett:} unclear --- no unique outcomes,
    so the framework is hard to apply. (Everettians
    typically \emph{claim} the theory is local --- but
    making this precise is notoriously difficult.)
  \item \textbf{Retrocausal models:} PI fails --- future settings
    influence past outcomes
  \end{itemize}

  \vspace{0.5em}
  \textbf{HH's focus:} Bell's theorem shows a fundamental
  limitation on \kw{local causal explanation} of quantum
  correlations. Does this mean we \emph{cannot understand} them?

  \vspace{0.3em}
  \small De Regt \& Dieks: understanding does not require causal
  explanation. Bohr: we need a new conceptual framework,
  not a new causal story.

\end{frame}

\begin{frame}{Taking stock --- and a question}

  \begin{itemize}
  \item QM achieves extraordinary predictive precision:
    the anomalous magnetic moment of the electron agrees
    with experiment to 12 significant figures
  \item Yet the theory is arguably \emph{inconsistent as stated}
    (the measurement problem), and after 100 years there is still
    no consensus on how to resolve it
  \item Each interpretation pays a serious price: non-locality,
    many worlds, anti-realism, new fundamental dynamics
  \end{itemize}

  \vspace{0.5em}
  Some physicists say: ``QM works, so who cares if we don't
  understand it?'' If we accept this, we have weakened our ground
  against:
  \begin{quote}\small
  ``AlphaFold predicts protein structures better than any human
  theory --- so who cares if we don't understand it?''
  \end{quote}

  \textbf{HH:} The goal of science is not mere prediction but
  understanding. QM is one thing; AI is another.

\end{frame}

%%% ============================================================
%%% PART II: AI AND SCIENTIFIC UNDERSTANDING
%%% ============================================================

\part{AI and Scientific Understanding}

\section{AI in Science}

\begin{frame}{}

  \begin{center}
  \includegraphics[width=0.78\textwidth]{boo.jpg}
  \end{center}

  \vspace{0.2em}
  \small\centering
  Humanities graduates at the University of Central Florida booed
  their commencement speaker after she called AI ``the next
  industrial revolution.'' \hfill \textit{New York Times}, May 2025

\end{frame}

\begin{frame}{Group Discussion}

  AlphaFold (2022) predicted the 3D structure of virtually every
  known protein --- a problem biologists had worked on for 50 years.
  AlphaFold is a deep neural network trained on known protein
  structures.

  \vspace{0.8em}
  \begin{enumerate}
  \item Did AlphaFold help us \textbf{understand} protein folding?
    Or did it only help us \textbf{predict} protein structures?
  \item Is there a meaningful difference between a neural network
    that predicts and a model that explains?
  \end{enumerate}

\end{frame}

\begin{frame}{Two perspectives on AI in science}

  \textbf{Theoretical}\enspace
  Science has always used tools to extend understanding.
  But AI tools seem borderline \emph{intelligent} themselves.
  Does that change what the task of science \emph{is}?

  \bigskip

  \textbf{Practical}\enspace Science produces two kinds of good:
  \textbf{knowledge} about the world, and \textbf{power} over the
  world. AI threatens to tip the balance away from knowledge.

  \bigskip What kind of science do we want to do? 

\end{frame}

\begin{frame}{A fact about AI: opacity}

  Deep Neural Networks (DNNs) --- including the transformer models
  behind AlphaFold and large language models --- share one crucial
  feature:

  \vspace{0.5em}
  \begin{center}
  \kw{Even the experts who build them do not know how they work}
  \end{center}

  \vspace{0.5em}
  \begin{itemize}
  \item There is a fundamental \kw{mismatch} between human reasoning
    and the path a DNN takes from inputs to outputs
  \item A DNN learns by adjusting billions of numerical weights
    through an optimization process; the resulting weights
    are not human-readable
  \item Joel Dudley (co-creator of the Deep Patient model):
    ``We can build these models, but we don't know how they work''
  \end{itemize}

\end{frame}

\begin{frame}{What a DNN looks like inside}

\begin{center}
\begin{tikzpicture}[scale=0.82]
  \foreach \i in {1,...,4} {
    \node[circle,draw,minimum size=0.55cm,fill=blue!10,font=\scriptsize]
      (in\i) at (0, {2.5-(\i-1)*0.88}) {};
  }
  \foreach \i in {1,...,5} {
    \node[circle,draw,minimum size=0.55cm,fill=gray!22,font=\scriptsize]
      (ha\i) at (2.2, {3.0-(\i-1)*0.82}) {};
  }
  \foreach \i in {1,...,5} {
    \node[circle,draw,minimum size=0.55cm,fill=gray!38,font=\scriptsize]
      (hb\i) at (4.4, {3.0-(\i-1)*0.82}) {};
  }
  \foreach \i in {1,2} {
    \node[circle,draw,minimum size=0.55cm,fill=red!12,font=\scriptsize]
      (out\i) at (6.6, {1.5-(\i-1)*1.0}) {};
  }
  \foreach \i in {1,...,4} {
    \foreach \j in {1,...,5} {
      \draw[gray!50,line width=0.25pt] (in\i) -- (ha\j);
    }
  }
  \foreach \i in {1,...,5} {
    \foreach \j in {1,...,5} {
      \draw[gray!50,line width=0.25pt] (ha\i) -- (hb\j);
    }
  }
  \foreach \i in {1,...,5} {
    \foreach \j in {1,2} {
      \draw[gray!50,line width=0.25pt] (hb\i) -- (out\j);
    }
  }
  \node[font=\scriptsize] at (0,3.2) {input};
  \node[font=\scriptsize] at (3.3,3.7) {hidden layers};
  \node[font=\scriptsize] at (6.6,2.2) {output};
  \node[font=\tiny,align=center,text=gray!65] at (3.3,-0.85)
    {each edge carries a weight $w_{ij}\in\mathbb{R}$,
     learned by gradient descent --- not designed to represent
     any physical quantity};
\end{tikzpicture}
\end{center}

\vspace{-0.2em}
\small The ``knowledge'' lives in the weights. These are the output of
gradient descent over training data --- not designed to represent any
physical mechanism or causal structure.

\end{frame}

\section{What Models Explain}

\begin{frame}{Models and understanding}

  Last week we discussed DN, causal-nomological, and
  unificationist accounts of explanation. Despite their
  differences, they share one assumption:

  \vspace{0.4em}
  \begin{center}
  \emph{A genuine explanation says something about what is
  going on in the world --- it does not merely predict.}
  \end{center}

  \vspace{0.4em}
  Contemporary philosophy of science stresses that explanations
  are typically underwritten by \kw{models}: representations
  with interpretable content.

  \vspace{0.4em}
  Such a model is not a black box for computing predictions.
  It has an \emph{inside} that scientists can inhabit,
  reason about, and transfer to new situations.

\end{frame}

\begin{frame}{Two examples}

  \textbf{Kinetic theory}\enspace gas $=$ molecules in random
  motion; explains temperature, pressure, and effusion rates
  by appealing to real physical processes.

  \bigskip

  \textbf{Schwarzschild metric}\enspace explains why light bends
  \emph{twice} the Newtonian prediction, because spatial
  curvature is a genuine feature of the geometry near a mass.

  \bigskip

  \small \textbf{Pointer to the literature:}
  van Fraassen, \textit{Scientific Representation: Paradoxes of
  Perspective} (Oxford, 2008) --- the central philosophical
  study of how models represent their target systems.

\end{frame}

\section{Sullivan's Argument}

\begin{frame}{Sullivan's question}

  \begin{quote}
  ``Are scientists trading understanding for some other epistemic
  or pragmatic good when they choose an opaque and complex machine
  learning model?''
  \end{quote}
  \hfill\small (Sullivan 2022, p.~110)

  \vspace{0.8em}
  \textbf{Discuss:} do you think scientists are making such a trade?
  Is that acceptable?

\end{frame}

\begin{frame}{The double meaning of ``model''}

  One terminological trap: the word \kw{model} does double duty.

  \vspace{0.6em}
  \textbf{Sense 1 --- scientific model}\par\smallskip
  Interpretable content: states a mechanism; has an inside that
  scientists can inhabit; satisfies CIT.
  \par{\small\emph{Examples: kinetic theory, Schwarzschild, BCS}}

  \vspace{0.6em}
  \textbf{Sense 2 --- ML model}\par\smallskip
  A parameterized function \emph{fitted to data by optimization}:
  weights set by gradient descent, not designed to represent anything.
  \par{\small\emph{Examples: AlphaFold, GPT, Deep Patient}}

\end{frame}

\begin{frame}{Why the ambiguity matters}

  The word ``model'' in ML descends from statistics, where it
  \emph{did} once have Sense~1 content (a Gaussian model, a linear
  model makes readable claims about a data-generating process).
  Neural networks inherited the term while quietly abandoning
  the interpretability.

  \vspace{0.8em}
  Sullivan's article uses ``model'' in both senses without clearly
  distinguishing them. She asks whether a \textbf{Sense~2} model
  can provide the understanding we normally expect from a
  \textbf{Sense~1} model.

  \vspace{0.5em}
  This is the right question --- but answering it requires
  acknowledging that there is a genuine gap to be bridged,
  not just a terminological overlap.

\end{frame}

\begin{frame}{Sullivan's main claim}

  Sullivan's answer to her own question: it is \emph{not}
  the complexity or opacity of a DNN that limits the understanding
  it can provide.

  \vspace{0.5em}
  The real culprit is \kw{link uncertainty}:
  \begin{quote}
  ``It is not implementation black-boxing that gets in the way of
  understanding, it is link uncertainty.''
  \end{quote}
  \hfill\small (Sullivan 2022, p.~116)

  \vspace{0.5em}
  \begin{itemize}
  \item Black-boxing implementation details does \emph{not}
    undermine understanding in general --- this is commonplace
    in science
  \item What blocks understanding is the absence of scientific
    evidence linking the model's outputs to real causal structure
    in the world
  \end{itemize}

  \vspace{0.4em}
  \textbf{Discuss:} do you think Sullivan is right about this?
  Can you give examples of acceptable and unacceptable black-boxing
  from physics?

\end{frame}

\begin{frame}{Black boxes}

  A \kw{black box}: a system whose inputs and outputs are known
  but whose internal workings are not accessible.

  \vspace{0.5em}
  \textbf{Implementation black-boxing} is ubiquitous in science
  --- and usually \emph{unproblematic}:
  \begin{itemize}
  \item A climate model's factorial subroutine may be iterative
    or recursive --- it makes no difference to the climatology
  \item Schelling's model gives the same result on a computer,
    a real board, or a Go board
  \end{itemize}

  \vspace{0.5em}
  \textbf{A familiar case:} many physicists treat QM itself as a
  \kw{quantum statistical algorithm} --- preparations go in,
  measurement statistics come out, do not ask what is happening
  inside. Part~I has been arguing this attitude is unsatisfactory.

\end{frame}

\begin{frame}{Black boxes in machine learning}

  DNNs are black boxes at a \emph{deeper} level: even the
  designers cannot identify which features the model has latched onto.

  \vspace{0.5em}
  The opacity goes beyond implementation details:
  \begin{itemize}
  \item Weights are set by optimization, not by design
  \item The modeller cannot predict which data features
    will be most salient
  \item Indirect probing (e.g.\ saliency maps) gives only
    approximate, high-level information
  \end{itemize}

  \vspace{0.5em}
  \textbf{Sullivan's claim:} even this deeper opacity does not
  block understanding. The real obstacle is something else.

\end{frame}

\begin{frame}{The Schelling analogy}

  Sullivan's central example: Schelling's \kw{checkerboard model}
  of residential segregation.
  \begin{itemize}
  \item Simple simulation: two types of agent, each moves if
    $>70\%$ of neighbours differ; result: a segregated board
  \item Provides \emph{how-possibly} understanding when the
    semantic mapping (coin = agent, move = choice) is read off
  \item Becomes \emph{how-actually} understanding only when
    empirical evidence connects the model to real populations
  \end{itemize}

  \vspace{0.5em}
  Sullivan's claim: the same structure applies to DNNs.
  Implementation opacity is not the obstacle --- link uncertainty is.

  \vspace{0.4em}
  \small \textbf{HH:} But notice the key difference. With Schelling,
  the semantic mapping was \emph{built in by the modeller from the
  start}. With Deep Patient, there is no such mapping. The very
  question of what the nodes and weights \emph{refer to} is
  unanswered --- and it is not obvious it is even askable.

\end{frame}

\begin{frame}{Sullivan's optimism and a challenge}

  \begin{quote}
  ``Once the link uncertainty is resolved, the deep patient model
  is able to explain and enable understanding of disease
  development.''
  \end{quote}
  \hfill\small (Sullivan 2022, p.~125)

  \vspace{0.8em}
  \textbf{HH:} How would the link uncertainty ever be resolved
  in this case?
  \begin{itemize}
  \item With Schelling: we look for evidence that individual
    preferences actually drive segregation --- this is a clearly
    statable empirical question
  \item With Deep Patient: what physical objects or processes
    would the learned weights be \emph{linked to}?
  \item The weights emerged from optimization over medical records.
    There is no design that maps them onto causal factors in disease.
    The link question cannot be stated in the same form.
  \end{itemize}

\end{frame}

\section{A Deeper Challenge}

\begin{frame}{The internal/external distinction}

  When a physicist derives a result, her brain is --- at the
  physical level --- just as opaque as a DNN.

  If you traced the neural mechanisms, you would not \emph{see}
  the reasoning. To understand why she gets the right answers,
  you have to get inside the \kw{structure of the reasoning itself}.

  \vspace{0.6em}
  \textbf{Leibniz's mill} (\emph{Monadologie}, \S{}17): even if
  the brain were enlarged to the size of a mill, and you could
  walk around inspecting all its machinery, you would never
  encounter \emph{thought} anywhere in the mechanism.

\end{frame}

\begin{frame}{Sullivan's blind spot?}

  Sullivan asks: how opaque is the implementation?

  \vspace{0.6em}
  But there is a prior question:

  \begin{center}
  \kw{Is there an internal perspective to be had at all?}
  \end{center}

  \vspace{0.5em}
  \begin{itemize}
  \item With the physicist: yes --- her reasoning has a logical
    structure you can inhabit
  \item With a DNN: \emph{open question}
  \end{itemize}

\end{frame}

\begin{frame}{Van Fraassen: a model must be a ``way the world could be''}

  Van Fraassen is the arch empiricist: he does not ask that our
  theories be \emph{true}, only \emph{empirically adequate}.

  \vspace{0.5em}
  Yet even van Fraassen would not be satisfied with Deep Patient
  as a final scientific result.

  \vspace{0.5em}
  His \kw{constructive empiricism} is a \emph{semantic} empiricism,
  not an eliminativist one. A scientific theory presents models ---
  structures interpretable as \emph{ways the world might be}.
  The scientist may be agnostic about truth, but she must be able
  to say: this is a \emph{possible description of reality}.

  \vspace{0.5em}
  A DNN does not present a way the world could be. It presents a
  function from inputs to outputs. The question of truth or
  falsity about disease mechanisms does not even arise.

\end{frame}

\begin{frame}{For reflection}

  \begin{enumerate}

  \item Should scientists rest content with having something as
    predictive as Deep Patient? Or does science have a further goal
    that remains unfulfilled?

  \item Does it matter whether we are talking about \emph{applied}
    science (predicting disease, protein structure) vs.
    \emph{fundamental} science (understanding the laws of nature)?

  \item We saw in Part~I that QM proceeds without consensus on
    what it means. Does this weaken or strengthen the case for
    demanding understanding from AI?

  \item Intimation: the problem of understanding another human
    being --- and why that seems different from understanding a DNN
    --- points toward questions that philosophy has been wrestling
    with for centuries.

  \end{enumerate}

\end{frame}

\nocite{*}

\begin{frame}[allowframebreaks]{Selected References}
\printbibliography[heading=none]
\end{frame}

\end{document}

%%% Local Variables:
%%% mode: latex
%%% TeX-master: t
%%% End:
